\documentclass[10pt,a4paper]{amsart}
\usepackage[margin=19mm]{geometry}
\usepackage{amsmath,amssymb,mathtools}
\usepackage[scr=rsfs,cal=euler]{mathalfa}
\usepackage{xfrac}
\usepackage[unicode,hidelinks,bookmarks=false]{hyperref}
\newcommand{\ie}{i.e.\ }
\newcommand{\cf}{cf.\ }
\newcommand{\resp}{respectively\ }
\newcommand{\Subsection}{\subsection}
\providecommand{\nobreakdash}{\nobreak\hbox{-}}
\setlength{\emergencystretch}{2em}
\multlinegap0pt
\usepackage{amssymb}
\usepackage{mathtools}
\usepackage{tikz}
\usepackage{float}
\newtheorem{lemma}{Lemma}
\newcommand{\Fpr}{\mathbb{P}}
\title{Asymptotics of the Meandering Number of a Cyclic Permutation}
\author{Charles Daly, Diaaeldin Taha}
\date{Source: arXiv:2606.12331v2 (CC BY 4.0)}
\begin{document}
\maketitle

\noindent\emph{Source and licence.} Asymptotics of the Meandering Number of a Cyclic Permutation. Version arXiv:2606.12331v2 (CC BY 4.0), distributed by arXiv under the Creative Commons Attribution 4.0 International licence, http://creativecommons.org/licenses/by/4.0/, as recorded in the arXiv licence metadata for this exact version and shown on the abstract page https://arxiv.org/abs/2606.12331v2. Full licence notice: \texttt{licenses/arxiv-2606.12331-CC-BY-4.0.txt}.\par\smallskip

\noindent\textit{Editorial scope.} Counted unit: the article's lemma lem:estimate\_prob (source0.tex lines 1026--1028). The article's complete original proof of it is delivered with it (source0.tex lines 1031--1068, one \texttt{proof} environment of 38 lines). No mathematics is written, repaired or re-derived: the statement and the proof are the article's own lines, in the article's own notation, and no retained line is substituted or re-flowed.  Retained uncounted as the source-local context the proof needs: lines 945--947; lines 951--957; lines 954--956; lines 998--1000.  Nothing was omitted from any retained span, so the package carries no omission record.  No retained line contains a citation, so the wrapper carries no bibliography and no bibliography entry.  No retained line contains a figure environment, an image-inclusion command or drawing code, so no image is delivered and the package declares no asset dependency.  Editorial formatting only, in addition: an AMS \texttt{amsart} preamble that restates the article's own declarations and macros used by the retained lines, namely the environments \texttt{lemma} on separate counters, together with the standard packages those lines need.  The retained environments print in this document as Lemma 1 in order of appearance, whereas the article prints the same items with the numbering of its own full document.  The complete native source of the article as distributed in the arXiv e-print for this exact version is bundled unchanged as \texttt{sources/202616/source0.tex}.
\begin{equation}\label{eq:binest}
\frac{1}{N+1}\exp\left(Nh(k/N)\right) \leq \binom{N}{k} \leq \exp\left(Nh(k/N)\right)
\end{equation}

\begin{lemma}\label{lem:binent}
There exists an $\eta > 0$ and an integer $L_{0}$, dependent on the choice of $\eta$, so that for all $L > L_{0}$, the inequalities of Equation \ref{eq:linequal} below hold. 
%%%%EQUATION_eq:linequal
\begin{equation}\label{eq:linequal}
4\left(\eta L + 3\right)^{2}\exp\left(2(L+1)h(3\eta)\right) \leq \exp\left(\frac{3}{4}h(1/3)L\right) \phantom{==} \log(L+1) \leq \frac{1}{16}h(1/3)L \phantom{==} \eta L \geq 4-3\eta
\end{equation} 
\end{lemma}

\begin{equation}
4\left(\eta L + 3\right)^{2}\exp\left(2(L+1)h(3\eta)\right) \leq \exp\left(\frac{3}{4}h(1/3)L\right) \phantom{==} \log(L+1) \leq \frac{1}{16}h(1/3)L \phantom{==} \eta L \geq 4-3\eta
\end{equation} 

\begin{lemma}[Separated Intervals Bounds]\label{lem:sepintbounds}
Let $r \leq s$ be non-negative integers and $L$ be a positive integer.  Pick some $\sigma$ randomly amongst the $L!$ linear orderings of $[L]$ under the uniform distribution.  Let $I_{s}(\sigma,r)$ and $I_{s}(\sigma_{0},r)$ denote the collections of $s$-elements subsets of $[L]$ which admit separated interval decompositions of at most $r$-intervals in the orderings $[L]_{\sigma}$ and $[L]_{0}$ respectively.  Define the random variable $X_{s,r}(\sigma) := |I_{s}(\sigma,r) \cap I_{s}(\sigma_{0},r)|$.  Then the probability that $X_{s,r}$ is greater than $0$ is bounded above $|I_{s}(\sigma_{0},r)|^{2}/\binom{L}{s}$.  Symbolically $\Fpr\left(X_{s,r} \geq 1 \right) \leq |I_{s}(\sigma_{0},r)|^{2}/\binom{L}{s}$.  
\end{lemma}

\begin{lemma}[Uniformly Bounded Probabilities]\label{lem:estimate_prob}
There exists a real $\eta > 0$ and positive integer $L_{0}$, which depends on $\eta$, such that the following holds.  We can find an $a > 0$ so that for any integer $L > L_{0}$, and all integers $s \in [L/3,2L/3]$, if we let $r = \lceil \eta L \rceil + 1$, then the probability that $X_{r,s}$ is bigger than $0$ is bounded above by $e^{-aL}$. 
\end{lemma}

\begin{proof}
Choose and $0 < \eta < 1/6$ satisfying $2h(3\eta) < \frac{1}{2}h(1/3)$ as in $\eta$ Lemma \ref{lem:binent}.  Now pick a positive integer $L_{0}$ as guaranteed by Lemma \ref{lem:binent} so that all three inequalities of Equation \ref{eq:linequal} hold for all $L > L_{0}$.  We now aim to bound the size of $I_{s}(\sigma_{0},r)$ for all $s \in [L/3, 2L/3]$. 
%%%%%%%
 Let $U$ be an $s$-element subset of $[L]$ expressible as a union of at most $r$ separated intervals in the order $[L]_{0}$.  Say there are $k$ such intervals.  Each interval $I_{i}$ may be assigned the pair of elements $(a_{i},b_{i})$ where $a_{i}$ is the smallest element of $I_{i}$, and $b_{i}$ is the immediate successor of the largest element of $I_{i}$.  To make this well defined, we define the immediate successor of $L-1$ to be $L \notin [L]$.  By definition $a_{i} < b_{i}$ for each $1 \leq i \leq k$.  Thus each such $U$ can be assigned the $k$-pairs of elements $\{(a_{i},b_{i})\}_{i=1}^{k}$ from $[L]\cup\{L\}$ which uniquely determine $U$, and its separated interval decomposition.  By hypothesis $k$ can be any number between $0$ and $r$, thus we have the bounds below.  
%%%%%%%EQUATION_eq:choose_est}
\begin{equation}\label{eq:choose_est}
|I_{s}(\sigma_{0},r)| \leq \sum_{k = 0}^{r} \binom{L+1}{2k} \leq \sum_{k=0}^{2r} \binom{L+1}{k}
\end{equation}
%%%%
%Recall that every bad set is expressible as a union of $r$-separated intervals in both $[L]$ and $[L]_{\sigma}$.  Consequently, our estimates show for a fixed $s \in [L/3,2L/3]$, the probability that a randomly selected permutation $\sigma$ admits a bad set of size $s$ is bounded above by $|I_{s}|^{2}/\binom{L}{s}$.  
%%%%
While this is a significant overestimate, it will be easier to work with than being precise.  In particular, we use the estimates in Equation \ref{eq:binest} to yield the inequality below for each $0 \leq k \leq 2r$.  
%%%%%
\begin{equation*}
\binom{L+1}{k} \leq \exp\left((L+1)h\left(\frac{k}{L+1}\right)\right)
\end{equation*}
By the third inequality of Equation \ref{eq:linequal}, $2r \leq 3\eta(L+1)$ and by definition of $\eta$, $3\eta < 1/2$.  Thus the monotonicity of $h$ on $[0,1/2]$ yields for each $0 \leq k \leq 2r$ that $h(k/(L+1)) \leq h(3\eta)$.  Combining these inequalities with Equation \ref{eq:choose_est}, we have for each $s \in [L/3,2L/3]$,
%%%%%%
\begin{equation*}
|I_{s}(\sigma_{0},r)| \leq \sum_{k=0}^{2r} \binom{L+1}{k} \leq (2r+1)\exp\left((L+1)h(3\eta)\right)
\end{equation*}
By squaring both sides of the above inequality, and using the fact that $L$ satisfies the first inequality of Equation \ref{eq:linequal}, we obtain for each $s \in [L/3,2L/3]$ the inequality below.
%%%%EQUATION_eq:numbound
\begin{equation}\label{eq:numbound}
|I_{s}(\sigma_{0},r)|^{2} \leq \exp\left(\frac{3}{4}h(1/3)L\right)
\end{equation}
Similarly, Equation \ref{eq:binest} and monotonicity of $h$ gives for each $s \in [L/3,2L/3]$ the inequality below.
%%%%%%%%EQUATION_eq:denbound}
\begin{equation}\label{eq:denbound}
\frac{1}{L+1}\exp\left(h(1/3)L\right) \leq \binom{L}{s}  
\end{equation}
Consequently for each $s \in [L/3,2L/3]$ we have that
%%%%%%%
\begin{equation*}
 \frac{|I_{s}(\sigma_{0},r)|^{2}}{\binom{L}{s}} \leq \left(L+1\right)\exp\left(-\frac{1}{4}h(1/3)L\right) \leq \exp\left(-\frac{3}{16}h(1/3)L\right)
\end{equation*}
The last inequality above is a consequence of the third inequality of Equation \ref{eq:linequal}.  Combining this result with Lemma \ref{lem:sepintbounds} establishes our result.
\end{proof}

\end{document}
